\chapter{Architectures and training}
\label{chap:archi}

This chapter is the engineering counterpart of \cref{chap:theory}.
\Cref{sec:archi-cnn,sec:archi-resnet,sec:archi-attn} describe the three
neural networks we trained, with their exact layer composition and the
intent behind each design choice. \Cref{sec:archi-compare} summarises
the three in a single comparison table. \Cref{sec:archi-train} then
documents the shared training pipeline.

\label{chap:archi-train}

The three models share the same \emph{interface}: input
$\mathbf{x} \in \R^{1 \times 512}$ (a single channel, a window of $512$
samples), output $\hat{p} \in \R^{7}$ (probabilities over the seven
classes). This makes them drop-in interchangeable for the training and
inference scripts, which select a model by name and otherwise behave
identically.

\section{Architecture A --- baseline 1D-CNN}
\label{sec:archi-cnn}

The first model, defined in
\texttt{architectures/saved\_models/1d\_cnn.py}, is a deliberately
compact convolutional network. It serves both as a baseline to beat
and as a sanity check that the data and the training pipeline are
correctly wired.

\paragraph{Structure.} The network is a linear stack of three
convolution--BatchNorm--ReLU--MaxPool blocks, followed by an adaptive
average pool and a small classifier head:

\begin{table}[H]
    \centering
    \footnotesize
    \begin{tabular}{@{}l p{0.55\linewidth} l@{}}
        \toprule
        \textbf{Block} & \textbf{Operation} & \textbf{Output shape} \\
        \midrule
        Input       & Raw window                                                                                       & $(1, 512)$ \\
        Block 1     & $\Conv(1{\to}16, k{=}7, p{=}3) \to \BN \to \ReLU \to \mathrm{MaxPool}(2)$                          & $(16, 256)$ \\
        Block 2     & $\Conv(16{\to}32, k{=}7, p{=}3) \to \BN \to \ReLU \to \mathrm{MaxPool}(2)$                         & $(32, 128)$ \\
        Block 3     & $\Conv(32{\to}64, k{=}5, p{=}2) \to \BN \to \ReLU \to \mathrm{MaxPool}(2)$                         & $(64, 64)$ \\
        Pool        & $\mathrm{AdaptiveAvgPool1d}(4)$                                                                  & $(64, 4)$ \\
        Classifier  & $\mathrm{Flatten} \to \mathrm{Dropout}(0.4) \to \mathrm{Linear}(256{\to}32) \to \ReLU \to \mathrm{Dropout}(0.2) \to \mathrm{Linear}(32{\to}7)$ & $(7)$ \\
        \bottomrule
    \end{tabular}
\end{table}

\paragraph{Design notes.}
\begin{itemize}[leftmargin=*, itemsep=0.2em]
    \item Kernel sizes $k=7,7,5$ are large enough to capture
    multi-sample shape primitives (ramps, peaks) at each scale, while
    staying well within the receptive-field budget.
    \item Channel widths $16 \to 32 \to 64$ double at each block,
    compensating the temporal compression by a factor of $2$ at every
    \texttt{MaxPool}.
    \item \texttt{AdaptiveAvgPool1d(4)} keeps four temporal bins after
    the convolutional trunk: this preserves a coarse notion of
    ``where in the window the activation peaks'', a useful piece of
    information for shapes such as \texttt{fast\_ascend} vs
    \texttt{fast\_descend} that mainly differ by the location of their
    peak.
    \item Two dropout layers ($p=0.4$ then $p=0.2$) in the
    fully-connected head act as a regulariser and prevent the small
    head from overfitting the convolutional features.
\end{itemize}

Total trainable parameters: $22\,727$.

\section{Architecture B --- 1D Residual Network}
\label{sec:archi-resnet}

The second model, defined in
\texttt{architectures/saved\_models/resNet.py}, replaces each
convolutional block of model A with a small \emph{residual block}, as
introduced in \cref{sec:resnet-theory}. It also replaces every
BatchNorm with GroupNorm (\cref{sec:gn-theory}).

\paragraph{Residual block.} A single block transforms an input of
channel count $c_{\text{in}}$ into an output of channel count
$c_{\text{out}}$ as
\begin{align*}
    \mathcal{F}(\mathbf{u}) &= \mathrm{GN} \!\circ\! \Conv_{k_3} \!\circ\! \ReLU \!\circ\! \mathrm{GN} \!\circ\! \Conv_{k_2} \!\circ\! \ReLU \!\circ\! \mathrm{GN} \!\circ\! \Conv_{k_1} (\mathbf{u}), \\
    \text{Block}(\mathbf{u}) &= \ReLU\bigl( \mathcal{F}(\mathbf{u}) + \mathrm{Shortcut}(\mathbf{u}) \bigr),
\end{align*}
where the shortcut is either the identity (when
$c_{\text{in}} = c_{\text{out}}$) or a $1{\times}1$ convolution
followed by GroupNorm. The kernel sizes
$(k_1, k_2, k_3)$ are $(15, 9, 5)$ in the first block and $(7, 5, 3)$
in the subsequent two; the first block uses a larger initial kernel to
capture broad context immediately.

\paragraph{Full structure.}

\begin{table}[H]
    \centering
    \footnotesize
    \begin{tabular}{@{}l p{0.55\linewidth} l@{}}
        \toprule
        \textbf{Block} & \textbf{Operation} & \textbf{Output shape} \\
        \midrule
        Input       & Raw window                                                              & $(1, 512)$ \\
        Block 1     & $\mathrm{Res}\bigl( c_{\text{in}}{=}1 \to c_{\text{out}}{=}16, \, k{=}(15,9,5) \bigr)$ & $(16, 512)$ \\
        Pool 1      & $\mathrm{MaxPool}(2)$                                                   & $(16, 256)$ \\
        Block 2     & $\mathrm{Res}\bigl( 16 \to 32, \, k{=}(7,5,3) \bigr)$                   & $(32, 256)$ \\
        Pool 2      & $\mathrm{MaxPool}(2)$                                                   & $(32, 128)$ \\
        Block 3     & $\mathrm{Res}\bigl( 32 \to 64, \, k{=}(7,5,3) \bigr)$                   & $(64, 128)$ \\
        Pool        & $\mathrm{AdaptiveAvgPool1d}(4)$                                          & $(64, 4)$ \\
        Classifier  & Same head as model~A                                                    & $(7)$ \\
        \bottomrule
    \end{tabular}
\end{table}

\paragraph{Design notes.}
\begin{itemize}[leftmargin=*, itemsep=0.2em]
    \item GroupNorm with $\min(8, c_{\text{out}})$ groups everywhere
    means train-time and inference-time behaviour are identical, which
    matters when deploying on real data with a different value
    distribution.
    \item The shortcut connections preserve the original signal
    through the convolutional trunk, allowing later blocks to focus on
    \emph{adding} discriminative information rather than rediscovering
    the input.
    \item The classifier head is intentionally identical to that of
    model A, so that the residual trunk is the only difference
    between the two networks --- a clean ablation.
\end{itemize}

Total trainable parameters: $74\,631$.

\section{Architecture C --- CNN+Attention}
\label{sec:archi-attn}

The third model, defined in
\texttt{architectures/saved\_models/cnn\_attention.py}, hybridises a
convolutional stem with a Transformer encoder
(\cref{sec:encoder-theory}). The motivation is to combine the
\emph{locality} of convolutions, which is well-suited to extracting
short shape primitives, with the \emph{global mixing} of attention,
which lets the network reason about how those primitives compose over
the whole window.

\paragraph{Convolutional stem.}
Three Conv--GroupNorm--ReLU blocks reduce the temporal length by a
factor of $8$ via stride-$2$ convolutions (so no explicit
\texttt{MaxPool} is needed):

\begin{table}[H]
    \centering
    \footnotesize
    \begin{tabular}{@{}l p{0.55\linewidth} l@{}}
        \toprule
        \textbf{Layer} & \textbf{Operation} & \textbf{Output shape} \\
        \midrule
        Input  & Raw window                                                                  & $(1, 512)$ \\
        Stem 1 & $\Conv(1{\to}16, k{=}7, s{=}2, p{=}3) \to \mathrm{GN}(8) \to \ReLU$           & $(16, 256)$ \\
        Stem 2 & $\Conv(16{\to}32, k{=}5, s{=}2, p{=}2) \to \mathrm{GN}(8) \to \ReLU$          & $(32, 128)$ \\
        Stem 3 & $\Conv(32{\to}64, k{=}3, s{=}2, p{=}1) \to \mathrm{GN}(8) \to \ReLU$          & $(64, 64)$ \\
        \bottomrule
    \end{tabular}
\end{table}

\paragraph{Tokenisation.}
The stem output is reshaped to a sequence of $T = 64$ tokens of
dimension $d = 64$. We prepend a learnable
$\texttt{[CLS]}$ token, giving $T' = 65$ tokens, and add a sinusoidal
positional encoding (\cref{eq:posenc}).

\paragraph{Transformer encoder.}
Two pre-LayerNorm encoder blocks (\cref{fig:transformer-block}) are
stacked. Each block has $h=4$ attention heads of dimension $d_h = 16$
and a feed-forward network with hidden size
$d_{\mathrm{ff}} = 4 \cdot d = 256$, GELU activation and dropout
$p=0.1$.

\paragraph{Classification.}
After the encoder, the $\texttt{[CLS]}$ row of the output is
LayerNorm-normalised and passed through a single linear layer
$\R^{d} \to \R^{7}$ to produce the class logits.

\paragraph{Design notes.}
\begin{itemize}[leftmargin=*, itemsep=0.2em]
    \item The stride-$2$ convolutional stem reduces the sequence length
    eightfold before attention is applied, which keeps the
    $O(T^2)$ cost of attention manageable at $65^2 \approx 4\,000$
    pairs.
    \item GroupNorm (not BatchNorm) in the stem makes the entire
    feature extraction amplitude-invariant on a per-window basis ---
    consistent with the window-mean centring of \cref{eq:window-norm}.
    \item The $\texttt{[CLS]}$ token is initialised with a small random
    vector (truncated normal, $\sigma = 0.02$) and is the only
    \emph{learnable} positional information; the sinusoidal
    encoding is shared and fixed.
\end{itemize}

Total trainable parameters: $109\,655$.

\section{Comparison of the three architectures}
\label{sec:archi-compare}

\Cref{tab:archi-compare} summarises the three architectures side by
side. The trend is clear: each step adds expressive power at the cost
of more parameters but \emph{not} much more training time, because the
deeper architectures still operate on heavily down-sampled features.

\begin{table}[H]
    \centering
    \footnotesize
    \begin{tabular}{lccc}
        \toprule
        & \textbf{1D-CNN} & \textbf{ResNet 1D} & \textbf{CNN+Attention} \\
        \midrule
        Convolutional trunk         & 3 plain blocks & 3 residual blocks & 3-layer stride-2 stem \\
        Normalisation               & BatchNorm      & GroupNorm(8)      & GroupNorm(8) \\
        Skip connections            & ---            & yes               & yes (inside Transformer) \\
        Long-range mixing           & ---            & ---               & Multi-head self-attention \\
        Output head                 & MLP(256\,$\to$\,32\,$\to$\,7) & MLP(256\,$\to$\,32\,$\to$\,7) & Linear(64\,$\to$\,7) on CLS \\
        Trainable parameters        & $22\,727$      & $74\,631$         & $109\,655$ \\
        Training time (35 epochs, RTX 5060) & $4'18''$ & $10'02''$         & $10'14''$ \\
        Inference throughput        & $10\,664$ win/s & $14\,232$ win/s   & $15\,184$ win/s \\
        \bottomrule
    \end{tabular}
    \caption{Side-by-side summary of the three architectures.
    Throughput is measured on the same RTX 5060 Laptop GPU at a stride
    of $10$ samples between successive windows.}
    \label{tab:archi-compare}
\end{table}

\section{Training pipeline}
\label{sec:archi-train}

The training pipeline is implemented once, in
\texttt{architectures/train.py}, and parameterised by the model name.

\subsection{Data preparation}

The raw signal is split temporally (80/20) into a train+val pool and a
test pool (\cref{sec:split-theory}), then converted into sliding
windows of size $512$ with stride $32$:
\begin{itemize}[leftmargin=*, itemsep=0.2em]
    \item Train+val pool: $249\,985$ windows
    (random 80/20 split: $199\,988$ train, $49\,997$ validation).
    \item Test pool: $62\,485$ windows.
\end{itemize}
Each window is centred individually and scaled by the global training
standard deviation $\sigma_{\text{train}} = 4.21$
(\cref{eq:window-norm}). $\sigma_{\text{train}}$ is saved in a JSON file
next to the model weights so that inference uses exactly the same
scaling.

\paragraph{Window class distribution.}
The labelling rule of \cref{sec:framing} yields roughly $31\%$ of
windows of class \texttt{Normal} and $11\%$--$12\%$ of each of the six
anomaly classes (\cref{tab:class-dist}).

\begin{table}[H]
    \centering
    \footnotesize
    \begin{tabular}{lrrr}
        \toprule
        Class & Count & Share & Class weight $w_c$ \\
        \midrule
        \texttt{Normal Background} & $77\,674$ & $31.1\%$ & $0.46$ \\
        \texttt{bell}              & $30\,218$ & $12.1\%$ & $1.18$ \\
        \texttt{bell\_sq}          & $28\,106$ & $11.2\%$ & $1.27$ \\
        \texttt{fast\_ascend}      & $27\,004$ & $10.8\%$ & $1.32$ \\
        \texttt{fast\_descend}     & $28\,769$ & $11.5\%$ & $1.24$ \\
        \texttt{M}                 & $30\,963$ & $12.4\%$ & $1.15$ \\
        \texttt{square}            & $27\,251$ & $10.9\%$ & $1.31$ \\
        \midrule
        Total                       & $249\,985$ & $100\%$ & --- \\
        \bottomrule
    \end{tabular}
    \caption{Window-level class distribution in the training set and
    the resulting class weights used in the cross-entropy loss.}
    \label{tab:class-dist}
\end{table}

\subsection{Optimisation}

All three models are trained with the recipe of
\cref{sec:train-theory}: Adam ($\eta=10^{-3}$, default
$\beta_1, \beta_2$), batch size $128$, weighted cross-entropy
\cref{eq:cross-entropy}, a $\texttt{ReduceLROnPlateau}$ scheduler with
patience $2$ and factor $0.5$, early stopping with patience $10$ and
minimum delta $5 \cdot 10^{-3}$, maximum $35$ epochs. No data
augmentation is applied --- the stochastic anomaly generation already
provides enough variability.

\subsection{Selection criterion}

The training script saves the weights with the lowest validation loss.
Validation F1 score (macro-averaged across the seven classes) is
tracked at every epoch for diagnostic purposes but is not used for
selection.
